\textbf{反射原理、投票问题与循环引理}

\textbf{约定}
走向右方的步记为 $R=(1,0)$，走向上方的步记为 $U=(0,1)$。
固定 $a$ 个 $R$、$b$ 个 $U$ 的路径总数为 $\binom{a+b}{b}$；
约定 $\binom{n}{k}=0$（$k<0$ 或 $k>n$）。

\textbf{反射原理：前缀不越界}
若 $a\ge b\ge0$，要求每个前缀满足 $\#U\le\#R$，则
\[
\#\text{合法路径}
=\binom{a+b}{b}-\binom{a+b}{b-1}
=\frac{a-b+1}{a+1}\binom{a+b}{b}.
\]
将首次到达 $\#U=\#R+1$ 之前的路径反射，可把违规路径与
$\binom{a+b}{b-1}$ 个无约束路径配对。
特别地，$a=b=n$ 时得到卡特兰数
$C_n=\frac1{n+1}\binom{2n}{n}$。

\textbf{带初始余量的上界}
若 $h\ge0$ 且 $b\le a+h$，要求每个前缀满足
$\#U\le\#R+h$，则
\[
\#\text{合法路径}
=\binom{a+b}{b}-\binom{a+b}{b-h-1}.
\]
终点若已违反 $b\le a+h$，合法路径数为 $0$；不能将上式直接套用。

\textbf{严格领先（Bertrand 投票问题）}
若 $a>b\ge0$，要求每个非空前缀都满足 $\#R>\#U$，则
\[
\#\text{合法路径}=\frac{a-b}{a+b}\binom{a+b}{b}.
\]
当 $a=b>0$ 时严格领先到终点不可能；$a=b=0$ 的空路径单独处理。

\textbf{循环引理}
把 $a$ 个 $+1$ 和 $b$ 个 $-1$ 排成环，且 $a>b$。
从 $a+b$ 个有标号起点逐个展开，恰有 $a-b$ 个起点使
每个非空前缀和都严格为正；即使环上出现周期，仍按起点位置计数。
适用于把“所有前缀合法”的线性序列转成循环移位计数。

\textbf{注意}
以上反射公式依赖单位步长和单侧直线边界；换成障碍点、
加权步长或多条边界时，须重新建立配对。
